\section{Working With Derivatives}

\subsection{Extremal Points}

\begin{definition} $ $
  \begin{itemize}
    \item $c \in D$ is called a {\bf local maximum point} of $f : D \rightarrow \mathbb{R}$ if there exists $\delta > 0$ such that for all $x \in D$ with $|x - c| < \delta$ we have
    \[f(x) \leq f(c).\]
\item $c \in D$ is called a {\bf local minimum point} of $f : D \rightarrow \mathbb{R}$ if there exists a $\delta > 0$ such that for all $x \in D$ with $|x - c| < \delta$ we have
    \[f(c) \leq f(x).\]
  \end{itemize}
\end{definition}

\begin{theorem}
    Let $D \subseteq \mathbb{R}$ and $f : D \rightarrow \mathbb{R}$. Suppose 
    \begin{itemize}
      \item $c$ is an interior point
      \item $c$ is either a local maximum or a local minimum
      \item f is differentiable at $c$.
    \end{itemize}
    Then \[f'(c) = 0.\]
\end{theorem}
\begin{proof}
    Suppose $c$ is a local maximum. Using this and that $c$ is an interior point we can find $\delta$ such that
    \[f(x) \leq f(c) \quad \forall \ x \in (c - \delta, c + \delta).\]
    Then for all $n > \frac{1}{\delta}$ we have
    \[f(c) \geq f(c - \frac{1}{n}), \quad f(c) \geq f(c + \frac{1}{n})\]
    so
    \[ \frac{f(c + \frac{1}{n}) - f(c)}{\frac{1}{n}} \leq 0 \leq \frac{f(c - \frac{1}{n}) - f(c)}{- \frac{1}{n}}.\]
    Taking the limit to infinity then gives that $f'(c)= 0$ by the sandwich theorem.
\end{proof}

\subsection{The Mean Value Theorem}

\begin{theorem} (Rolle's Theorem)

  \noindent Let $a < b$ and $f : [a , b] \rightarrow \mathbb{R}$. Suppose $f$ is continuous on $[a,b]$ and differentiable on $(a, b)$. If $f(a) = f(b)$ then there exists $c \in (a, b)$ with $f'(c) = 0$.
\end{theorem}
\begin{proof}
  By the Weierstrass sextreme value theorem $f$ has a global minimum $c_1 \in [a,b]$ and a global maximum $c_2 \in [a,b]$. If either of these are in $(a,b)$ we are done. Otherwise the minimum and maximum are equal so for any $x$,
  \[f(a) \leq f(x) \leq f(x) \ \ \implies \ \ f(x) = f(a)\]
  and $f'(x) = 0$.
\end{proof}

\begin{theorem} (Mean Value Theorem)
    
  \noindent Let $a < b$ and $f : [a,b] \rightarrow \mathbb{R}$ be continuous on $[a,b]$ and differentiable on $(a,b)$. Then there exists $c \in (a,b)$ such that
  \[f'(c) = \frac{f(b) - f(a)}{b - a}.\]
\end{theorem}
\begin{proof}
    Let $g(x) := f(x) - x \frac{f(b) - f(a)}{b - a}$. Then
    \begin{align*}
      g(b) - g(a) &= f(b) - f(a) - b \frac{f(b) - f(a)}{b -a} + a \frac{f(b) - f(a)}{b -a}\\
       &= f(b) - f(a) - \frac{(b - a)(f(b) - f(a))}{b - a} \\
      &= 0
    .\end{align*}
    Thus we can use Rolle's theorem to find $c$ with $g'(c) = 0$, which can be rearanged for the desired result.
\end{proof}

\begin{corollary}
  Let $a < b$ and $f: [a,b] \rightarrow \mathbb{R}$ be continuous on $[a,b]$ and differentiable on $(a,b)$. If $f'(x) = 0$ for all $x \in (a,b)$ then $f$ is constant on $[a,b]$.
\end{corollary}
\begin{proof}
  For $x \in (a, b]$, by the mean value theorem there exists $c \in (a, x)$ such that
  \[f'(c) = 0 =  \frac{f(x) - f(a)}{x - a} \quad \implies f(x) = f(a).\]
\end{proof}

\begin{theorem}
  Let $a < b$, and $f,g : [a,b] \rightarrow \mathbb{R}$ be continuous on $[a,b]$ and differentiable on $(a,b)$. Then there exists $c \in (a,b)$ such that
  \[(f(b) -f(a))(g'(c)) = (g(b) - g(a))f'(c) \]
  or alternatively that
  \[\frac{f'(c)}{g'(c)} = \frac{f(c) - f(a)}{g(b) - g(a)}.\]
\end{theorem}
\begin{proof}
  Consider $h: [a,b] \rightarrow \mathbb{R}$ with
  \[h(x) = ( f(b) - f(a))g(x) - (g(b) - g(a))f(x).\]
  Then it can be shown that
  \[h(a) = f(b) g(a) - f(a) g(b) = h(b).\]
  Then as $h$ is differentiable we can use Rolle's theorem to find a point c where $h'(c) = 0$ which can be rearanged for the desired result.
\end{proof}
